\section{Síntesis y ampliación}

Lo que más vale conservar de esta clase no es la lista de modelos, sino un marco de análisis de tres niveles. El primer nivel es \textbf{representation/objective}: qué información conserva cada uno de language token, continuous vision feature, discrete image code y noisy latent. El segundo nivel es \textbf{architecture/systems}: cómo projection, experts, cross attention, ZeRO, parallelism y long context llevan la representación a un hardware en el que se puede entrenar. El tercer nivel es \textbf{data/evaluation}: cómo instruction, preference, caption, GUI trace, video coverage y benchmark definen el comportamiento que el modelo termina aprendiendo.

\subsection{Unificación de los tres ejes}

\begin{importantbox}{Panorama unificado final del curso}
\begin{enumerate}
  \item \textbf{El LLM aporta la base secuencial de uso general: }scaling, alignment y systems hacen del modelo grande una interfaz reutilizable.
  \item \textbf{La multimodality expone el cuello de botella de la interfaz: }los detalles visuales, la alta resolución, el espacio bidimensional y la estructura temporal del video obligan al modelo a cambiar el bridge, el objective y la compression.
  \item \textbf{Los data determinan los límites del comportamiento: }la architecture puede ampliar el rango que el modelo es capaz de ajustar, pero benchmark, annotation, recaption, verification y failure coverage determinan lo que el sistema aprende en realidad.
\end{enumerate}
\end{importantbox}

\begin{warningbox}{Los tres errores de generalización más frecuentes}
No conviertas «la misma loss de preentrenamiento» en un teorema de «la misma capacidad»; no conviertas una sola GUI trace en agent reliability; no conviertas la practical advantage de diffusion en una prueba de imposibilidad de la autoregression. Los tres errores provienen de borrar las diferencias entre los niveles de evidencia presentados en clase.
\end{warningbox}

\subsection{Checklist para reproducir experimentos}

Para reproducir cualquiera de los modelos de esta clase, se debe publicar al mismo tiempo, como mínimo, architecture version, data mixture/provenance, resolution/tokenización, training compute, inference steps, prompt/evaluation protocol y repeated-trial failures. Publicar solo el headline score impide distinguir entre data, algorithm y architecture.

\begin{lstlisting}[caption={Lista unificada de reproducción para LLM/VLM/diffusion}]
pin(code_commit, model_checkpoint, tokenizer_or_vae)
document(data_sources, filters, synthetic_fraction, licenses)
report(train_flops, gpu_hours, parallelism, peak_memory)
report(input_resolution, context_length, sampling_steps)
freeze(prompts, decoding, benchmark_version, test_split)
publish(ablations, confidence_intervals, failures, costs)
\end{lstlisting}

\subsection{Preguntas de autoevaluación}

\begin{multicols}{2}
\footnotesize
\begin{enumerate}
  \setlength{\itemsep}{0pt}
  \setlength{\parskip}{0pt}
  \item ¿En qué difiere la información de condicionamiento de GLM blank infilling respecto de un MLM puro y de una autoregression pura?
  \item ¿Por qué la scaling law es un asignador de presupuesto y no una garantía de rendimiento?
  \item ¿Por qué la perplexity no puede compararse directamente entre tokenizers distintos?
  \item ¿Qué tipo de inference memory reduce principalmente GQA?
  \item ¿A qué aplica sharding cada uno de ZeRO-1/2/3?
  \item ¿En qué se diferencia activation checkpointing de model checkpointing?
  \item ¿Qué dimensión divide cada uno de TP, PP y context parallelism?
  \item ¿Por qué el principal costo de interacción de long context puede estar en el prefill?
  \item DPO no tiene un reward model independiente; ¿por qué sigue dependiendo de reward assumptions?
  \item ¿Cómo muestra el ejemplo de CogQA que data/algorithm/architecture son intercambiables?
  \item ¿Qué dos tipos de trabajo realiza a la vez el Q-Former de BLIP-2?
  \item La projection de LLaVA es sencilla, pero ¿por qué el OCR todavía puede ser deficiente?
  \item ¿Qué capacidad buscan proteger los vision experts de CogVLM?
  \item ¿Por qué CogAgent necesita una ruta lateral de alta resolución?
  \item ¿Qué capacidades no puede demostrar una sola web trace?
  \item ¿En qué difieren Vary y GLM-4V en el tratamiento del OCR bottleneck?
  \item ¿Qué ventajas y desventajas aporta el image tokenizer a la autoregressive generation?
  \item ¿Por qué diffusion suele aprovechar mejor la GPU que la AR decoding con lote de tamaño uno?
  \item ¿Por qué Relay Diffusion necesita cambiar la noise correlation entre resolutions?
  \item ¿En qué difieren las rutas de condition de adaLN y de cross attention?
  \item En una Sora-like recipe, ¿qué es un problema de representación y qué es un problema de systems/data?
  \item ¿Por qué los annotated video pairs no garantizan que se aprendan las physical rules?
\end{enumerate}
\end{multicols}

\subsection{Lecturas adicionales}

\begin{multicols}{2}
\footnotesize
\begin{itemize}
  \setlength{\itemsep}{0.15em}
  \setlength{\parskip}{0pt}
  \item Du et al., \href{https://arxiv.org/abs/2103.10360}{GLM: General Language Model Pretraining with Autoregressive Blank Infilling}: para entender la unificación de las funciones objetivo.
  \item Rajbhandari et al., \href{https://arxiv.org/abs/1910.02054}{ZeRO}; Shoeybi et al., \href{https://arxiv.org/abs/1909.08053}{Megatron-LM}: para entender la partición del estado y el paralelismo de modelo.
  \item Rafailov et al., \href{https://arxiv.org/abs/2305.18290}{Direct Preference Optimization}: para entender el pairwise preference objective.
  \item Li et al., \href{https://arxiv.org/abs/2301.12597}{BLIP-2}; Liu et al., \href{https://arxiv.org/abs/2304.08485}{LLaVA}: para comparar el Q-Former con el projection bridge.
  \item Wang et al., \href{https://arxiv.org/abs/2311.03079}{CogVLM}; Hong et al., \href{https://arxiv.org/abs/2312.08914}{CogAgent}: para entender los vision experts y el GUI high-resolution path.
  \item Ho et al., \href{https://arxiv.org/abs/2006.11239}{DDPM}; Peebles and Xie, \href{https://arxiv.org/abs/2212.09748}{DiT}; Esser et al., \href{https://arxiv.org/abs/2403.03206}{Stable Diffusion 3}: del diffusion objective al Transformer denoiser.
  \item Snell et al., \href{https://arxiv.org/abs/2409.01236}{Same Pre-training Loss, Better Downstream}: lectura post-lecture que sirve de contraejemplo a la afirmación categórica de la clase de que «la loss determina la capacidad».
\end{itemize}
\end{multicols}

\end{document}
